\begin{comment}
Inside each section, communicate the following:
* The scientific evolution chronologically
    * Conceptual feasibility
    * Practical experiments
    * Validated technology
* Identify limitations
* Highlight unresolved problems
\end{comment}

The literature review conducted for this paper can be divided into three themes:
\begin{itemize}
    \item GNSS-R Small-Satellite Missions
    \item Formation Flying and Formation Maintenance for Small Satellites
    \item Integrated Satellite Mission Simulation and Lifetime Analysis    
\end{itemize}


\section{GNSS-R Small-Satellite Missions}

\begin{itemize}
    \item Scientific evolution:
    \begin{enumerate}
        \item GNSS as a signal of opportunity
        \item GNSS-R feasibility for different remote sensing applications. Bistatic radar principle
        \item Early demonstrations of GNSS-R remote sensing on aircraft mission
        \item Spacecraft demonstrations
        \item Validation of technology for different applications
        \item Expansion of GNSS-R Mission applications (include jamming/spoofing detection)
    \end{enumerate}
    \item Mapping useful GNSS-R missions after validation of technology at spacecraft altitudes  
    \item Feasibility of cooperative swarms and their respective formation types w. benefits
\end{itemize}

Reflected \gls{gnss} signals were first discussed as a



\section{Formation Flying and Formation Maintenance for Small Satellites}

\subsection{6U satellite formation-control lifetime}
"Published nanosatellite and 6U CubeSat formation-flying missions commonly plan formation-keeping demonstration phases on the order of several months, with six months being a representative value. Some missions have spacecraft design lifetimes of one year or more, but the duration over which active formation-keeping can be maintained depends strongly on the propulsion system, control strategy, orbit altitude, separation distance, and required relative-position accuracy."


For the VISORS mission, a formation-flying pair of 6U CubeSats in LEO, the published concept of operations states a six-month mission lifetime; this should be interpreted as a mission-specific design lifetime rather than a general expected lifetime for formation-maintaining CubeSats in LEO.\cite{lightsey2024visors_conops}

(3U satellite) CPOD launched on May 25, 2022, and NASA reported on June 23, 2023 that the mission had ended after fuel depletion, implying fuel depletion occurred no later than about 13 months after launch. \cite{nasa2023cpod}

Formation maintenance and formation-flying experiments consumed a significant part of the propellant budget, but in this mission they did not appear to determine end-of-life. The spacecraft completed the formation-flying objectives with substantial remaining propellant.\cite{roth2016canx45_flight_results}


\subsection{Deployment speed}
The NanoRacks NRCSD IDD states that a CubeSat must be compatible with a deployment/ejection velocity of 0.5 to 2.0 m/s from the NanoRacks CubeSat Deployer. It also gives a related attitude-rate requirement: the CubeSat must withstand up to 5 deg/s/axis tip-off rate, while noting that the NRCSD target tip-off rate is less than 2 deg/s/axis. \cite{nanoracks2018nrcsd_idd}


The NanoRacks External CubeSat Deployer interface requirements specify that CubeSats shall tolerate deployment velocities of 0.5-2.5 m/s and tip-off rates up to 5 deg/s/axis at ejection from the NRCSD-E. \cite{nanoracks2018nrcsde_idd}


ISISPACE specifies that its QuadPack CubeSat deployer provides a release speed of 0.8-1.8 m/s. Also does support 6U deployment \cite{isispace2026quadpack}

The Dhruva Space DSOD-6U deployer supports 6U CubeSat deployment and specifies an ejection velocity below 2 m/s. Also support 2kg per U \cite{satsearch2026dhruva_dsod6u}


\section{Integrated Satellite Mission Simulation and Lifetime Analysis}

\section{Research gaps and thesis contributions}



